\section{相关工作}\label{ux76f8ux5173ux5de5ux4f5c}

\letteredsubsection{A. 多标签图像识别}\label{a.ux591aux6807ux7b7eux56feux50cfux8bc6ux522b}

不同于单标签分类任务，多标签图像识别允许一幅图像同时对应多个非互斥类别，因此更需要处理标签共现、标签依赖和类别不平衡等问题\hyperref[_Ref213949872]{{[}14{]}}。早期方法如Binary Relevance\hyperref[_Ref213949896]{{[}15{]}}通常将多标签任务简单拆分为多个独立二分类问题，该类方法容易忽略标签之间的相关关系，在复杂视觉场景中，往往会导致预测结果不完整或矛盾。

为缓解上述问题，近年来的多标签图像识别方法逐渐关注标签依赖建模和样本不平衡处理。ML-GCN\hyperref[_Ref213949382]{{[}10{]}}根据标签共现统计构建标签图，并通过图卷积网络传播标签关联信息；Asymmetric Loss (ASL)\hyperref[_Ref213949927]{{[}16{]}}对正负样本引入不同的聚焦参数，并通过负样本概率移位机制抑制易分类负样本的损失贡献，从而缓解多标签分类中的正负样本不平衡问题；Query2Label\hyperref[_Ref213950246]{{[}17{]}}进一步利用Transformer查询机制建模标签交互，并在VOC\hyperref[_Ref213950282]{{[}18{]}}等自然图像数据集上取得良好效果。这些工作表明，将标签关系显式或隐式地融入模型，有助于提升多标签预测的完整性和准确性。

在医疗眼底图像中，多标签识别具有更强的临床意义，因为同一只眼可能同时存在多种疾病，忽略疾病共存关系容易造成误检和漏检\hyperref[_Ref213949362]{{[}9{]}}。已有研究从不同角度探索眼底多疾病识别问题。例如，Cheng等\hyperref[_Ref213950349]{{[}19{]}}将ML-GCN迁移至眼底病灶多标签分类任务，通过构建眼底病灶语料和GloVe标签表示建模病灶标签依赖；DCNet\hyperref[_Ref213950338]{{[}20{]}}通过空间相关模块捕捉左右眼眼底图像特征之间的像素级关联，并融合双眼特征进行患者级多标签眼科疾病识别；DKCNet\hyperref[_Ref213950343]{{[}21{]}}利用注意力模块生成判别性区域特征，并结合通道重标定与重采样策略缓解眼科数据集中的类别不平衡问题。

总体而言，现有方法已经能够在一定程度上建模标签依赖或缓解类别不平衡，但在眼底多疾病共存场景中，疾病间关系复杂且病灶表现具有局部性，仅依赖标签共现统计或通用注意力机制仍可能难以稳定刻画疾病标签与病灶特征之间的对应关系。因此，面向眼底图像构建更自适应的标签-特征关联仍具有研究意义。

\letteredsubsection{B. 域泛化}\label{b.ux57dfux9002ux5e94ux57dfux6cdbux5316ux4e0eux5355ux57dfux6cdbux5316}

在计算机视觉任务中，模型在训练域上获得较好性能后，部署到新机构或新设备采集的数据上时，常会受到成像条件、光照环境和数据分布差异的影响而性能下降。这种域偏移在医疗成像中尤为突出，不同医院、设备和采集流程之间的数据差异会直接影响模型泛化能力\hyperref[_Ref213949439]{{[}22{]}}。域适应（Domain Adaptation, DA）是缓解域偏移的常见思路，其核心是在训练过程中利用源域和目标域数据进行分布对齐。例如，基于最大均值差异的Deep Domain Confusion\hyperref[_Ref213949448]{{[}23{]}}和相关性对齐方法CORAL\hyperref[_Ref213949453]{{[}24{]}}通过缩小特征分布差异提升迁移性能；DANN\hyperref[_Ref213949479]{{[}25{]}}采用对抗训练学习域不变特征；CycleGAN\hyperref[_Ref213949490]{{[}26{]}}及相关无监督域适应方法\hyperref[_Ref213949530]{{[}27{]}}则通过图像风格转换或跨域对比学习缓解域间差距。然而，DA通常要求训练阶段能够访问目标域数据，这一前提在医疗数据隐私敏感或目标机构数据不可提前获取的场景中并不总是成立。

与DA不同，域泛化（Domain Generalization, DG）旨在不访问目标域数据的条件下学习可迁移表示，并泛化到未知目标域\hyperref[_Ref213949554]{{[}28{]}}。典型DG方法包括MixStyle\hyperref[_Ref213949562]{{[}29{]}}，通过混合实例风格统计模拟潜在域偏移；DomainBed\hyperref[_Ref213949607]{{[}30{]}}提供统一基准以评估不同DG算法；FACT\hyperref[_Ref213949788]{{[}31{]}}从频域角度分解域特定和域不变信息；RSC\hyperref[_Ref213949796]{{[}32{]}}通过抑制主导特征迫使模型关注更多与类别相关的线索；基于元学习的DG方法\hyperref[_Ref213949803]{{[}33{]}}则通过模拟训练域与测试域差异来提升泛化能力。多数DG研究通常利用多个带标注源域学习跨域稳定表示，这有助于模型捕捉更丰富的域间变化，但也对标签空间统一、数据处理流程协调和跨机构协作提出较高要求。然而，在真实医疗场景中，上述条件并不总是能够满足，导致多机构数据往往难以集中或同步使用。即使通过联邦学习等方式避免集中存储，也会引入通信、隐私机制、非独立同分布数据和协同优化等额外复杂度。

在此背景下，单域泛化（Single Domain Generalization, SDG）作为更受限的泛化设置，仅使用一个源域学习面向未知域的可迁移特征。医疗图像中的single-source DG研究\hyperref[_Ref213949863]{{[}34{]}}也说明，在无法集中使用多机构数据时，从单一源域中挖掘更稳定的疾病相关表示具有现实价值。因此，仅利用单一源域学习可迁移的疾病相关表示，是更贴近跨机构眼底疾病识别部署约束的研究方向。

\letteredsubsection{C. 隐私保护}\label{c.ux9690ux79c1ux4fddux62a4}

医疗图像和电子健康记录包含高度敏感的个人信息，模型训练、参数共享或云端推理过程均可能带来隐私泄露风险\hyperref[_Ref213950358]{{[}35{]}}。在医学人工智能应用中，隐私保护不仅关系到患者权益，也受到数据合规要求的约束。常见的隐私保护方案包括联邦学习、加密计算、安全多方计算和差分隐私等。其中，联邦学习（Federated Learning, FL）允许多个客户端在本地训练模型并上传参数或梯度，从而减少原始数据集中共享的需求\hyperref[_Ref213950376]{{[}36{]}}。例如，FedProx\hyperref[_Ref213950391]{{[}37{]}}通过添加近端项来处理客户端异构数据，缓解了协同优化不稳定问题。

除联邦学习外，为进一步防范隐私泄露，加密机制可以提供更强隐私保障。同态加密（Homomorphic Encryption, HE）允许在密文上执行计算\hyperref[_Ref213950427]{{[}38{]}}，并已被探索用于医疗影像分析\hyperref[_Ref213950433]{{[}39{]}}，但通常伴随较高计算开销；安全多方计算（Secure Multi-Party Computation, SMPC）\hyperref[_Ref213950438]{{[}40{]}}可在多方协作时保护各方输入隐私；可信执行环境（Trusted Execution Environment, TEE）\hyperref[_Ref213950453]{{[}41{]}}则通过硬件隔离提供更安全的计算环境。差分隐私（Differential Privacy, DP）通过向查询结果或训练梯度中注入噪声，为单个样本的信息泄露风险提供可量化约束\hyperref[_Ref213950443]{{[}42{]}}。其中，高斯差分隐私\hyperref[_Ref213950449]{{[}43{]}}通过高斯噪声刻画隐私与效用之间的折衷，在深度学习梯度保护中具有较强适用性。

现有隐私保护机制可以降低医疗数据泄露风险，但在跨域或泛化任务中仍面临额外挑战。联邦学习、同态加密和安全多方计算往往带来通信、计算或系统协同成本；差分隐私虽然部署形式相对直接，但噪声注入可能对模型精度造成影响，且这种影响在不同类别或不同数据分布上可能并不均衡\hyperref[_Ref213950458]{{[}44{]}}，这一定程度上增加了多域泛化条件下的隐私保护难度。因此，在单域泛化训练流程中探索更可控的隐私--性能折衷，对于隐私受限场景下的跨机构眼底疾病识别具有实际意义。
